\begin{PSbox}[unbreakable]{obj}{Learning Objectives}

After completing this module, students will be able to:

\begin{enumerate}
\def\labelenumi{\arabic{enumi}.}
\tightlist
\item
  State and explain the Law of Large Numbers
\item
  State and apply the Central Limit Theorem
\item
  Explain why Gaussian distributions appear so frequently in engineering
\item
  Use CLT for approximations
\item
  Demonstrate CLT visually using MATLAB Monte Carlo simulations
\end{enumerate}

\end{PSbox}

\PSsection{11.1}{Introduction: Why Is the Gaussian Everywhere?}

The Gaussian distribution appears in:

\begin{itemize}
\tightlist
\item
  Thermal noise, measurement errors, manufacturing tolerances, aggregate interference, stock fluctuations
\end{itemize}

\textbf{The reason:} All these quantities are sums (or averages) of many small, independent random effects. The CLT guarantees convergence to Gaussian regardless of the underlying distribution.

\PSsection{11.2}{Law of Large Numbers (LLN)}

\PSsubsection{Weak Law of Large Numbers}

For X\textsubscript{1}, X\textsubscript{2}, …, X\textsubscript{n} i.i.d. with mean \ensuremath{\mu} and finite variance \ensuremath{\sigma}\textsuperscript{2}:

\PSdisplay{\bar{X}_n = \frac{1}{n}\sum_{i=1}^n X_i \xrightarrow{P} \mu \quad \text{as } n \to \infty}

(``Convergence in probability'': P(\textbar{}\ensuremath{\bar{X}}\textsubscript{n} - \ensuremath{\mu}\textbar{} \textgreater{} \ensuremath{\varepsilon}) \ensuremath{\to} 0 for any \ensuremath{\varepsilon} \textgreater{} 0)

\PSsubsection{Strong Law of Large Numbers}

\PSdisplay{\bar{X}_n \xrightarrow{a.s.} \mu \quad \text{as } n \to \infty}

(``Almost sure convergence'': P(\ensuremath{\bar{X}}\textsubscript{n} \ensuremath{\to} \ensuremath{\mu}) = 1)

\PSsubsection{Engineering Significance}

\begin{itemize}
\tightlist
\item
  \textbf{Sample average converges to true mean} as sample size grows
\item
  Justifies estimating means from data
\item
  Foundation of Monte Carlo simulation (simulate many times \ensuremath{\to} average converges)
\item
  Justifies: ``If I measure enough times, my average will be close to the truth''
\end{itemize}

\PSsubsection{How Fast?}

By Chebyshev: P(\textbar{}\ensuremath{\bar{X}}\textsubscript{n} - \ensuremath{\mu}\textbar{} \textgreater{} \ensuremath{\varepsilon}) \ensuremath{\le} \ensuremath{\sigma}\textsuperscript{2}/(n\ensuremath{\varepsilon}\textsuperscript{2})

For 95\% confidence that error \textless{} \ensuremath{\varepsilon}: need n \ensuremath{\ge} \ensuremath{\sigma}\textsuperscript{2}/(0.05 \ensuremath{\cdot} \ensuremath{\varepsilon}\textsuperscript{2})

\PSsection{11.3}{Central Limit Theorem (CLT)}

\PSsubsection{Formal Statement}

For X\textsubscript{1}, X\textsubscript{2}, …, X\textsubscript{n} i.i.d. with mean \ensuremath{\mu} and finite variance \ensuremath{\sigma}\textsuperscript{2}:

\PSdisplay{Z_n = \frac{\bar{X}_n - \mu}{\sigma/\sqrt{n}} = \frac{\sum_{i=1}^n X_i - n\mu}{\sigma\sqrt{n}} \xrightarrow{d} N(0,1)}

\PSsubsection{Practical Form}

For large n:\PSdisplay{\bar{X}_n \approx N\left(\mu, \frac{\sigma^2}{n}\right)}

\PSdisplay{S_n = \sum_{i=1}^n X_i \approx N(n\mu, n\sigma^2)}

\PSsubsection{Conditions}

\begin{enumerate}
\def\labelenumi{\arabic{enumi}.}
\tightlist
\item
  X\textsubscript{1}, …, X\textsubscript{n} are independent
\item
  Identically distributed (can be relaxed with Lyapunov/Lindeberg CLT)
\item
  Finite mean \ensuremath{\mu} and finite variance \ensuremath{\sigma}\textsuperscript{2}
\item
  n is ``large enough'' (rule of thumb: n \ensuremath{\ge} 30, but depends on distribution shape)
\end{enumerate}

\PSsubsection{What CLT Does NOT Require}

\begin{itemize}
\tightlist
\item
  Does NOT require X\textsubscript{i} to be Gaussian
\item
  Does NOT require X\textsubscript{i} to be continuous
\item
  Does NOT require any specific distribution
\end{itemize}

\PSsection{11.4}{Rate of Convergence}

How fast does the CLT ``kick in''?

{\def\LTcaptype{none} % do not increment counter
\begin{longtable}[c]{@{}cc@{}}
\toprule\noalign{}
\rowcolor{PSnavy}\PSth{Original Distribution} & \PSth{n needed for good approximation} \\
\midrule\noalign{}
\endhead
\bottomrule\noalign{}
\endlastfoot
Symmetric (uniform) & n \ensuremath{\approx} 5-10 \\
Mildly skewed & n \ensuremath{\approx} 20-30 \\
Highly skewed (exponential) & n \ensuremath{\approx} 30-50 \\
Extremely skewed & n \ensuremath{\approx} 50-100+ \\
\end{longtable}
}

\textbf{Rule:} More symmetric distributions \ensuremath{\to} faster convergence.

\PSsection{11.5}{Visual Demonstration}

\PSsubsection{Sum of Uniform RVs \ensuremath{\to} Gaussian}

\begin{itemize}
\tightlist
\item
  n=1: Uniform (flat)
\item
  n=2: Triangular
\item
  n=3: Starts looking bell-shaped
\item
  n=5: Nearly Gaussian
\item
  n=12: Practically indistinguishable from Gaussian
\end{itemize}

\PSsubsection{Sum of Exponential RVs \ensuremath{\to} Gaussian}

\begin{itemize}
\tightlist
\item
  n=1: Exponential (heavily right-skewed)
\item
  n=3: Still skewed but less so (Gamma(3))
\item
  n=10: Nearly symmetric, bell-shaped
\item
  n=30: Essentially Gaussian
\end{itemize}

\PSsubsection{Binomial \ensuremath{\to} Normal Approximation}

Binomial(n, p) \ensuremath{\approx} N(np, np(1-p)) for large n

This is CLT applied to sum of Bernoulli trials!

\begin{itemize}
\tightlist
\item
  Continuity correction: P(X \ensuremath{\le} k) \ensuremath{\approx} \ensuremath{\Phi}((k + 0.5 - np)/(\ensuremath{\surd}(np(1-p))))
\end{itemize}

\PSsection{11.6}{Engineering Significance}

\PSsubsection{Why Thermal Noise is Gaussian}

Thermal noise = aggregate effect of billions of electrons moving randomly. Each electron’s contribution is tiny and independent \ensuremath{\to} CLT applies \ensuremath{\to} noise is Gaussian.

\PSsubsection{Why Aggregate Interference is Gaussian}

In a cellular network, total interference = sum of signals from many independent users. CLT \ensuremath{\to} interference \ensuremath{\approx} Gaussian (for many users).

\PSsubsection{Why Measurement Errors are Gaussian}

Total error = sum of many independent small error sources (calibration, quantization, thermal, vibration, etc.) \ensuremath{\to} CLT \ensuremath{\to} error \ensuremath{\approx} Gaussian.

\PSsubsection{Why Manufacturing Variations are Gaussian}

Resistor value = nominal + sum of many small process variations (doping, etching, temperature, etc.) \ensuremath{\to} CLT \ensuremath{\to} value \ensuremath{\approx} Gaussian around nominal.

\PSsection{11.7}{MATLAB Monte Carlo Demonstrations}

\begin{PSbox}[unbreakable]{ex}{Example 1: CLT with Uniform Distribution}

\tcbset{PScodemode/.style={unbreakable}}

\begin{Shaded}
\begin{Highlighting}[]
\CommentTok{\%\% CLT: Sum of Uniform RVs approaches Gaussian}
\VariableTok{N\_experiments} \OperatorTok{=} \FloatTok{50000}\OperatorTok{;}
\VariableTok{n\_values} \OperatorTok{=}\NormalTok{ [}\FloatTok{1}\OperatorTok{,} \FloatTok{2}\OperatorTok{,} \FloatTok{5}\OperatorTok{,} \FloatTok{12}\OperatorTok{,} \FloatTok{30}\NormalTok{]}\OperatorTok{;}

\VariableTok{figure}\OperatorTok{;}
\KeywordTok{for} \VariableTok{idx} \OperatorTok{=} \FloatTok{1}\OperatorTok{:}\VariableTok{length}\NormalTok{(}\VariableTok{n\_values}\NormalTok{)}
    \VariableTok{n} \OperatorTok{=} \VariableTok{n\_values}\NormalTok{(}\VariableTok{idx}\NormalTok{)}\OperatorTok{;}
    
    \CommentTok{\% Generate sum of n Uniform(0,1) RVs}
    \VariableTok{X} \OperatorTok{=} \VariableTok{sum}\NormalTok{(}\VariableTok{rand}\NormalTok{(}\VariableTok{n}\OperatorTok{,} \VariableTok{N\_experiments}\NormalTok{)}\OperatorTok{,} \FloatTok{1}\NormalTok{)}\OperatorTok{;}
    
    \CommentTok{\% Standardize: Z = (X {-} nμ)/(σ√n) where μ=0.5, σ²=1/12}
    \VariableTok{mu\_sum} \OperatorTok{=} \VariableTok{n} \OperatorTok{*} \FloatTok{0.5}\OperatorTok{;}
    \VariableTok{sigma\_sum} \OperatorTok{=} \VariableTok{sqrt}\NormalTok{(}\VariableTok{n} \OperatorTok{/} \FloatTok{12}\NormalTok{)}\OperatorTok{;}
    \VariableTok{Z} \OperatorTok{=}\NormalTok{ (}\VariableTok{X} \OperatorTok{{-}} \VariableTok{mu\_sum}\NormalTok{) }\OperatorTok{/} \VariableTok{sigma\_sum}\OperatorTok{;}
    
    \VariableTok{subplot}\NormalTok{(}\FloatTok{2}\OperatorTok{,} \FloatTok{3}\OperatorTok{,} \VariableTok{idx}\NormalTok{)}\OperatorTok{;}
    \VariableTok{histogram}\NormalTok{(}\VariableTok{Z}\OperatorTok{,} \FloatTok{60}\OperatorTok{,} \SpecialStringTok{\textquotesingle{}Normalization\textquotesingle{}}\OperatorTok{,} \SpecialStringTok{\textquotesingle{}pdf\textquotesingle{}}\NormalTok{)}\OperatorTok{;}
    \VariableTok{hold} \VariableTok{on}\OperatorTok{;}
    \VariableTok{z} \OperatorTok{=} \VariableTok{linspace}\NormalTok{(}\OperatorTok{{-}}\FloatTok{4}\OperatorTok{,} \FloatTok{4}\OperatorTok{,} \FloatTok{200}\NormalTok{)}\OperatorTok{;}
    \VariableTok{plot}\NormalTok{(}\VariableTok{z}\OperatorTok{,} \VariableTok{normpdf}\NormalTok{(}\VariableTok{z}\NormalTok{)}\OperatorTok{,} \SpecialStringTok{\textquotesingle{}r\textquotesingle{}}\OperatorTok{,} \SpecialStringTok{\textquotesingle{}LineWidth\textquotesingle{}}\OperatorTok{,} \FloatTok{2}\NormalTok{)}\OperatorTok{;}
    \VariableTok{title}\NormalTok{(}\VariableTok{sprintf}\NormalTok{(}\SpecialStringTok{\textquotesingle{}n = \%d\textquotesingle{}}\OperatorTok{,} \VariableTok{n}\NormalTok{))}\OperatorTok{;}
    \VariableTok{xlabel}\NormalTok{(}\SpecialStringTok{\textquotesingle{}Standardized Sum\textquotesingle{}}\NormalTok{)}\OperatorTok{;} \VariableTok{xlim}\NormalTok{([}\OperatorTok{{-}}\FloatTok{4} \FloatTok{4}\NormalTok{])}\OperatorTok{;}
    \KeywordTok{if} \VariableTok{idx} \OperatorTok{==} \FloatTok{1}\OperatorTok{,} \VariableTok{ylabel}\NormalTok{(}\SpecialStringTok{\textquotesingle{}PDF\textquotesingle{}}\NormalTok{)}\OperatorTok{;} \KeywordTok{end}
\KeywordTok{end}
\VariableTok{sgtitle}\NormalTok{(}\SpecialStringTok{\textquotesingle{}CLT: Sum of Uniform RVs → Gaussian\textquotesingle{}}\NormalTok{)}\OperatorTok{;}
\end{Highlighting}
\end{Shaded}

\end{PSbox}

\begin{PSbox}[unbreakable]{ex}{Example 2: CLT with Exponential Distribution}

\tcbset{PScodemode/.style={unbreakable}}

\begin{Shaded}
\begin{Highlighting}[]
\CommentTok{\%\% CLT: Sample Mean of Exponential(1) converges to Gaussian}
\VariableTok{lambda} \OperatorTok{=} \FloatTok{1}\OperatorTok{;}  \VariableTok{mu} \OperatorTok{=} \FloatTok{1}\OperatorTok{;}  \VariableTok{sigma} \OperatorTok{=} \FloatTok{1}\OperatorTok{;}
\VariableTok{N\_experiments} \OperatorTok{=} \FloatTok{100000}\OperatorTok{;}
\VariableTok{n\_values} \OperatorTok{=}\NormalTok{ [}\FloatTok{1}\OperatorTok{,} \FloatTok{3}\OperatorTok{,} \FloatTok{10}\OperatorTok{,} \FloatTok{30}\OperatorTok{,} \FloatTok{100}\NormalTok{]}\OperatorTok{;}

\VariableTok{figure}\OperatorTok{;}
\KeywordTok{for} \VariableTok{idx} \OperatorTok{=} \FloatTok{1}\OperatorTok{:}\VariableTok{length}\NormalTok{(}\VariableTok{n\_values}\NormalTok{)}
    \VariableTok{n} \OperatorTok{=} \VariableTok{n\_values}\NormalTok{(}\VariableTok{idx}\NormalTok{)}\OperatorTok{;}
    
    \CommentTok{\% Generate sample means}
    \VariableTok{X} \OperatorTok{=} \VariableTok{exprnd}\NormalTok{(}\FloatTok{1}\OperatorTok{,} \VariableTok{n}\OperatorTok{,} \VariableTok{N\_experiments}\NormalTok{)}\OperatorTok{;}
    \VariableTok{X\_bar} \OperatorTok{=} \VariableTok{mean}\NormalTok{(}\VariableTok{X}\OperatorTok{,} \FloatTok{1}\NormalTok{)}\OperatorTok{;}
    
    \VariableTok{subplot}\NormalTok{(}\FloatTok{2}\OperatorTok{,} \FloatTok{3}\OperatorTok{,} \VariableTok{idx}\NormalTok{)}\OperatorTok{;}
    \VariableTok{histogram}\NormalTok{(}\VariableTok{X\_bar}\OperatorTok{,} \FloatTok{80}\OperatorTok{,} \SpecialStringTok{\textquotesingle{}Normalization\textquotesingle{}}\OperatorTok{,} \SpecialStringTok{\textquotesingle{}pdf\textquotesingle{}}\NormalTok{)}\OperatorTok{;}
    \VariableTok{hold} \VariableTok{on}\OperatorTok{;}
    \VariableTok{x} \OperatorTok{=} \VariableTok{linspace}\NormalTok{(}\FloatTok{0}\OperatorTok{,} \VariableTok{max}\NormalTok{(}\VariableTok{X\_bar}\NormalTok{)}\OperatorTok{,} \FloatTok{200}\NormalTok{)}\OperatorTok{;}
    \VariableTok{plot}\NormalTok{(}\VariableTok{x}\OperatorTok{,} \VariableTok{normpdf}\NormalTok{(}\VariableTok{x}\OperatorTok{,} \VariableTok{mu}\OperatorTok{,} \VariableTok{sigma}\OperatorTok{/}\VariableTok{sqrt}\NormalTok{(}\VariableTok{n}\NormalTok{))}\OperatorTok{,} \SpecialStringTok{\textquotesingle{}r\textquotesingle{}}\OperatorTok{,} \SpecialStringTok{\textquotesingle{}LineWidth\textquotesingle{}}\OperatorTok{,} \FloatTok{2}\NormalTok{)}\OperatorTok{;}
    \VariableTok{title}\NormalTok{(}\VariableTok{sprintf}\NormalTok{(}\SpecialStringTok{\textquotesingle{}n = \%d\textquotesingle{}}\OperatorTok{,} \VariableTok{n}\NormalTok{))}\OperatorTok{;}
    \VariableTok{xlabel}\NormalTok{(}\SpecialStringTok{\textquotesingle{}Sample Mean\textquotesingle{}}\NormalTok{)}\OperatorTok{;}
\KeywordTok{end}
\VariableTok{sgtitle}\NormalTok{(}\SpecialStringTok{\textquotesingle{}CLT: Mean of Exponential → Gaussian\textquotesingle{}}\NormalTok{)}\OperatorTok{;}
\end{Highlighting}
\end{Shaded}

\end{PSbox}

\begin{PSbox}[unbreakable]{ex}{Example 3: Binomial Normal Approximation}

\tcbset{PScodemode/.style={unbreakable}}

\begin{Shaded}
\begin{Highlighting}[]
\CommentTok{\%\% Binomial(n,p) approximated by Gaussian}
\VariableTok{n} \OperatorTok{=} \FloatTok{50}\OperatorTok{;} \VariableTok{p} \OperatorTok{=} \FloatTok{0.3}\OperatorTok{;}
\VariableTok{k} \OperatorTok{=} \FloatTok{0}\OperatorTok{:}\VariableTok{n}\OperatorTok{;}
\VariableTok{pmf\_binom} \OperatorTok{=} \VariableTok{binopdf}\NormalTok{(}\VariableTok{k}\OperatorTok{,} \VariableTok{n}\OperatorTok{,} \VariableTok{p}\NormalTok{)}\OperatorTok{;}

\CommentTok{\% Normal approximation}
\VariableTok{mu\_approx} \OperatorTok{=} \VariableTok{n}\OperatorTok{*}\VariableTok{p}\OperatorTok{;}
\VariableTok{sigma\_approx} \OperatorTok{=} \VariableTok{sqrt}\NormalTok{(}\VariableTok{n}\OperatorTok{*}\VariableTok{p}\OperatorTok{*}\NormalTok{(}\FloatTok{1}\OperatorTok{{-}}\VariableTok{p}\NormalTok{))}\OperatorTok{;}
\VariableTok{pdf\_normal} \OperatorTok{=} \VariableTok{normpdf}\NormalTok{(}\VariableTok{k}\OperatorTok{,} \VariableTok{mu\_approx}\OperatorTok{,} \VariableTok{sigma\_approx}\NormalTok{)}\OperatorTok{;}

\VariableTok{figure}\OperatorTok{;}
\VariableTok{bar}\NormalTok{(}\VariableTok{k}\OperatorTok{,} \VariableTok{pmf\_binom}\OperatorTok{,} \SpecialStringTok{\textquotesingle{}FaceAlpha\textquotesingle{}}\OperatorTok{,} \FloatTok{0.5}\NormalTok{)}\OperatorTok{;}
\VariableTok{hold} \VariableTok{on}\OperatorTok{;}
\VariableTok{plot}\NormalTok{(}\VariableTok{k}\OperatorTok{,} \VariableTok{pdf\_normal}\OperatorTok{,} \SpecialStringTok{\textquotesingle{}r{-}\textquotesingle{}}\OperatorTok{,} \SpecialStringTok{\textquotesingle{}LineWidth\textquotesingle{}}\OperatorTok{,} \FloatTok{2}\NormalTok{)}\OperatorTok{;}
\VariableTok{xlabel}\NormalTok{(}\SpecialStringTok{\textquotesingle{}k\textquotesingle{}}\NormalTok{)}\OperatorTok{;} \VariableTok{ylabel}\NormalTok{(}\SpecialStringTok{\textquotesingle{}Probability\textquotesingle{}}\NormalTok{)}\OperatorTok{;}
\VariableTok{title}\NormalTok{(}\VariableTok{sprintf}\NormalTok{(}\SpecialStringTok{\textquotesingle{}Binomial(\%d, \%.1f) vs Normal(\%.1f, \%.2f)\textquotesingle{}}\OperatorTok{,} \VariableTok{n}\OperatorTok{,} \VariableTok{p}\OperatorTok{,} \VariableTok{mu\_approx}\OperatorTok{,} \VariableTok{sigma\_approx}\OperatorTok{\^{}}\FloatTok{2}\NormalTok{))}\OperatorTok{;}
\VariableTok{legend}\NormalTok{(}\SpecialStringTok{\textquotesingle{}Binomial (exact)\textquotesingle{}}\OperatorTok{,} \SpecialStringTok{\textquotesingle{}Normal approximation\textquotesingle{}}\NormalTok{)}\OperatorTok{;}
\end{Highlighting}
\end{Shaded}

\end{PSbox}

\PSsubsection{Example 4: Complete Monte Carlo CLT Demonstration}

\tcbset{PScodemode/.style={unbreakable}}

\begin{Shaded}
\begin{Highlighting}[]
\CommentTok{\%\% Full CLT demonstration: non{-}Gaussian → Gaussian via averaging}
\CommentTok{\% Generate N samples from a VERY non{-}Gaussian distribution}
\CommentTok{\% (mixture: 80\% from Exp(1), 20\% from Exp(0.1))}
\VariableTok{N\_experiments} \OperatorTok{=} \FloatTok{100000}\OperatorTok{;}
\VariableTok{n\_samples} \OperatorTok{=} \FloatTok{50}\OperatorTok{;}

\CommentTok{\% Generate raw data}
\VariableTok{raw} \OperatorTok{=} \VariableTok{zeros}\NormalTok{(}\VariableTok{n\_samples}\OperatorTok{,} \VariableTok{N\_experiments}\NormalTok{)}\OperatorTok{;}
\KeywordTok{for} \VariableTok{i} \OperatorTok{=} \FloatTok{1}\OperatorTok{:}\VariableTok{N\_experiments}
    \VariableTok{mask} \OperatorTok{=} \VariableTok{rand}\NormalTok{(}\VariableTok{n\_samples}\OperatorTok{,} \FloatTok{1}\NormalTok{) }\OperatorTok{\textless{}} \FloatTok{0.8}\OperatorTok{;}
    \VariableTok{raw}\NormalTok{(}\VariableTok{mask}\OperatorTok{,} \VariableTok{i}\NormalTok{) }\OperatorTok{=} \VariableTok{exprnd}\NormalTok{(}\FloatTok{1}\OperatorTok{,} \VariableTok{sum}\NormalTok{(}\VariableTok{mask}\NormalTok{)}\OperatorTok{,} \FloatTok{1}\NormalTok{)}\OperatorTok{;}
    \VariableTok{raw}\NormalTok{(}\OperatorTok{\textasciitilde{}}\VariableTok{mask}\OperatorTok{,} \VariableTok{i}\NormalTok{) }\OperatorTok{=} \VariableTok{exprnd}\NormalTok{(}\FloatTok{10}\OperatorTok{,} \VariableTok{sum}\NormalTok{(}\OperatorTok{\textasciitilde{}}\VariableTok{mask}\NormalTok{)}\OperatorTok{,} \FloatTok{1}\NormalTok{)}\OperatorTok{;}
\KeywordTok{end}

\CommentTok{\% True mean and variance of mixture}
\VariableTok{mu\_true} \OperatorTok{=} \FloatTok{0.8}\OperatorTok{*}\FloatTok{1} \OperatorTok{+} \FloatTok{0.2}\OperatorTok{*}\FloatTok{10}\OperatorTok{;}  \CommentTok{\% = 2.8}
\VariableTok{var\_true} \OperatorTok{=} \FloatTok{0.8}\OperatorTok{*}\NormalTok{(}\FloatTok{1}\OperatorTok{+}\FloatTok{1}\OperatorTok{\^{}}\FloatTok{2}\NormalTok{) }\OperatorTok{+} \FloatTok{0.2}\OperatorTok{*}\NormalTok{(}\FloatTok{100}\OperatorTok{+}\FloatTok{10}\OperatorTok{\^{}}\FloatTok{2}\NormalTok{) }\OperatorTok{{-}} \VariableTok{mu\_true}\OperatorTok{\^{}}\FloatTok{2}\OperatorTok{;}  \CommentTok{\% Second moment {-} mean\^{}2}

\CommentTok{\% Compute sample means}
\VariableTok{X\_bar} \OperatorTok{=} \VariableTok{mean}\NormalTok{(}\VariableTok{raw}\OperatorTok{,} \FloatTok{1}\NormalTok{)}\OperatorTok{;}

\VariableTok{figure}\OperatorTok{;}
\VariableTok{subplot}\NormalTok{(}\FloatTok{2}\OperatorTok{,}\FloatTok{1}\OperatorTok{,}\FloatTok{1}\NormalTok{)}\OperatorTok{;}
\VariableTok{histogram}\NormalTok{(}\VariableTok{raw}\NormalTok{(}\OperatorTok{:,}\FloatTok{1}\NormalTok{)}\OperatorTok{,} \FloatTok{100}\OperatorTok{,} \SpecialStringTok{\textquotesingle{}Normalization\textquotesingle{}}\OperatorTok{,} \SpecialStringTok{\textquotesingle{}pdf\textquotesingle{}}\NormalTok{)}\OperatorTok{;}
\VariableTok{title}\NormalTok{(}\SpecialStringTok{\textquotesingle{}Original Distribution (highly skewed mixture)\textquotesingle{}}\NormalTok{)}\OperatorTok{;}
\VariableTok{xlabel}\NormalTok{(}\SpecialStringTok{\textquotesingle{}X\textquotesingle{}}\NormalTok{)}\OperatorTok{;} \VariableTok{ylabel}\NormalTok{(}\SpecialStringTok{\textquotesingle{}PDF\textquotesingle{}}\NormalTok{)}\OperatorTok{;}

\VariableTok{subplot}\NormalTok{(}\FloatTok{2}\OperatorTok{,}\FloatTok{1}\OperatorTok{,}\FloatTok{2}\NormalTok{)}\OperatorTok{;}
\VariableTok{histogram}\NormalTok{(}\VariableTok{X\_bar}\OperatorTok{,} \FloatTok{100}\OperatorTok{,} \SpecialStringTok{\textquotesingle{}Normalization\textquotesingle{}}\OperatorTok{,} \SpecialStringTok{\textquotesingle{}pdf\textquotesingle{}}\NormalTok{)}\OperatorTok{;}
\VariableTok{hold} \VariableTok{on}\OperatorTok{;}
\VariableTok{x} \OperatorTok{=} \VariableTok{linspace}\NormalTok{(}\VariableTok{min}\NormalTok{(}\VariableTok{X\_bar}\NormalTok{)}\OperatorTok{,} \VariableTok{max}\NormalTok{(}\VariableTok{X\_bar}\NormalTok{)}\OperatorTok{,} \FloatTok{200}\NormalTok{)}\OperatorTok{;}
\VariableTok{plot}\NormalTok{(}\VariableTok{x}\OperatorTok{,} \VariableTok{normpdf}\NormalTok{(}\VariableTok{x}\OperatorTok{,} \VariableTok{mu\_true}\OperatorTok{,} \VariableTok{sqrt}\NormalTok{(}\VariableTok{var\_true}\OperatorTok{/}\VariableTok{n\_samples}\NormalTok{))}\OperatorTok{,} \SpecialStringTok{\textquotesingle{}r\textquotesingle{}}\OperatorTok{,} \SpecialStringTok{\textquotesingle{}LineWidth\textquotesingle{}}\OperatorTok{,} \FloatTok{2}\NormalTok{)}\OperatorTok{;}
\VariableTok{title}\NormalTok{(}\VariableTok{sprintf}\NormalTok{(}\SpecialStringTok{\textquotesingle{}Sample Mean (n=\%d) — Gaussian by CLT!\textquotesingle{}}\OperatorTok{,} \VariableTok{n\_samples}\NormalTok{))}\OperatorTok{;}
\VariableTok{xlabel}\NormalTok{(}\SpecialStringTok{\textquotesingle{}Sample Mean\textquotesingle{}}\NormalTok{)}\OperatorTok{;} \VariableTok{ylabel}\NormalTok{(}\SpecialStringTok{\textquotesingle{}PDF\textquotesingle{}}\NormalTok{)}\OperatorTok{;}
\VariableTok{legend}\NormalTok{(}\SpecialStringTok{\textquotesingle{}Empirical\textquotesingle{}}\OperatorTok{,} \SpecialStringTok{\textquotesingle{}Gaussian approximation\textquotesingle{}}\NormalTok{)}\OperatorTok{;}
\end{Highlighting}
\end{Shaded}

\PSsection{11.8}{Approximations Using CLT}

\PSsubsection{General Procedure}

To approximate P(S\textsubscript{n} \ensuremath{\le} x) where S\textsubscript{n} = X\textsubscript{1} + … + X\textsubscript{n}:

\begin{enumerate}
\def\labelenumi{\arabic{enumi}.}
\tightlist
\item
  Compute \ensuremath{\mu} = E[X\textsubscript{i}], \ensuremath{\sigma}\textsuperscript{2} = Var(X\textsubscript{i})
\item
  S\textsubscript{n} \ensuremath{\approx} N(n\ensuremath{\mu}, n\ensuremath{\sigma}\textsuperscript{2})
\item
  P(S\textsubscript{n} \ensuremath{\le} x) \ensuremath{\approx} \ensuremath{\Phi}((x - n\ensuremath{\mu})/(\ensuremath{\sigma}\ensuremath{\surd}n))
\end{enumerate}

\PSsubsection{Continuity Correction (for Discrete RVs)}

When approximating a discrete sum by Gaussian:

\begin{itemize}
\tightlist
\item
  P(X \ensuremath{\le} k) \ensuremath{\approx} \ensuremath{\Phi}((k + 0.5 - n\ensuremath{\mu})/(\ensuremath{\sigma}\ensuremath{\surd}n))
\item
  P(X = k) \ensuremath{\approx} \ensuremath{\Phi}((k + 0.5 - n\ensuremath{\mu})/(\ensuremath{\sigma}\ensuremath{\surd}n)) - \ensuremath{\Phi}((k - 0.5 - n\ensuremath{\mu})/(\ensuremath{\sigma}\ensuremath{\surd}n))
\end{itemize}

\begin{PSbox}[unbreakable]{ex}{Example: Packet Errors}

1000 packets, each error probability 0.02. X = total errors.\PSbr{}Exact: X \textasciitilde{} Binomial(1000, 0.02)\PSbr{}CLT approx: X \ensuremath{\approx} N(20, 19.6)

P(X \textgreater{} 30) \ensuremath{\approx} 1 - \ensuremath{\Phi}((30.5 - 20)/\ensuremath{\surd}19.6) = 1 - \ensuremath{\Phi}(2.37) = 0.0089

\end{PSbox}

\PSsection{11.9}{Practice Problems}

\begin{PSbox}[unbreakable]{prob}{Problem 1}

A factory produces bolts. Length X \textasciitilde{} distribution with \ensuremath{\mu} = 10cm, \ensuremath{\sigma} = 0.2cm. A batch of 100 bolts is measured.

\begin{enumerate}
\def\labelenumi{(\alph{enumi})}
\tightlist
\item
  What is the approximate distribution of \ensuremath{\bar{X}}? (b) P(\ensuremath{\bar{X}} \textgreater{} 10.05)? (c) P(9.96 \textless{} \ensuremath{\bar{X}} \textless{} 10.04)?
\end{enumerate}

\PSsolution{} 

\begin{enumerate}
\def\labelenumi{(\alph{enumi})}
\tightlist
\item
  \ensuremath{\bar{X}} \textasciitilde{} N(10, 0.04/100) = N(10, 0.0004), \ensuremath{\sigma}\textsubscript{X}̄ = 0.02
\item
  P(\ensuremath{\bar{X}} \textgreater{} 10.05) = P(Z \textgreater{} 2.5) = 0.0062
\item
  P(9.96 \textless{} \ensuremath{\bar{X}} \textless{} 10.04) = \ensuremath{\Phi}(2) - \ensuremath{\Phi}(-2) = 0.9545
\end{enumerate}

\end{PSbox}

\begin{PSbox}[unbreakable]{prob}{Problem 2}

A communication link has 10000 bits with BER = 0.001. Approximate P(more than 15 errors) using CLT.

\PSsolution{} X \textasciitilde{} Binomial(10000, 0.001). \ensuremath{\mu}=10, \ensuremath{\sigma}\textsuperscript{2}=9.99.\PSbr{}P(X \textgreater{} 15) \ensuremath{\approx} 1 - \ensuremath{\Phi}((15.5-10)/\ensuremath{\surd}9.99) = 1 - \ensuremath{\Phi}(1.74) = 0.0409

\end{PSbox}

\begin{PSbox}[unbreakable]{prob}{Problem 3}

Write MATLAB code to verify the CLT for chi-squared(1) distribution (very skewed) and show convergence for n=2,5,10,50.

\PSsolution{} 

\tcbset{PScodemode/.style={unbreakable}}

\begin{Shaded}
\begin{Highlighting}[]
\VariableTok{N} \OperatorTok{=} \FloatTok{100000}\OperatorTok{;} \VariableTok{ns} \OperatorTok{=}\NormalTok{ [}\FloatTok{2} \FloatTok{5} \FloatTok{10} \FloatTok{50}\NormalTok{]}\OperatorTok{;}
\VariableTok{figure}\OperatorTok{;}
\KeywordTok{for} \VariableTok{i} \OperatorTok{=} \FloatTok{1}\OperatorTok{:}\FloatTok{4}
    \VariableTok{X\_bar} \OperatorTok{=} \VariableTok{mean}\NormalTok{(}\VariableTok{chi2rnd}\NormalTok{(}\FloatTok{1}\OperatorTok{,} \VariableTok{ns}\NormalTok{(}\VariableTok{i}\NormalTok{)}\OperatorTok{,} \VariableTok{N}\NormalTok{)}\OperatorTok{,} \FloatTok{1}\NormalTok{)}\OperatorTok{;}
    \VariableTok{subplot}\NormalTok{(}\FloatTok{2}\OperatorTok{,}\FloatTok{2}\OperatorTok{,}\VariableTok{i}\NormalTok{)}\OperatorTok{;}
    \VariableTok{histogram}\NormalTok{(}\VariableTok{X\_bar}\OperatorTok{,} \FloatTok{80}\OperatorTok{,} \SpecialStringTok{\textquotesingle{}Normalization\textquotesingle{}}\OperatorTok{,} \SpecialStringTok{\textquotesingle{}pdf\textquotesingle{}}\NormalTok{)}\OperatorTok{;} \VariableTok{hold} \VariableTok{on}\OperatorTok{;}
    \VariableTok{x} \OperatorTok{=} \VariableTok{linspace}\NormalTok{(}\FloatTok{0}\OperatorTok{,} \FloatTok{3}\OperatorTok{,} \FloatTok{200}\NormalTok{)}\OperatorTok{;}
    \VariableTok{plot}\NormalTok{(}\VariableTok{x}\OperatorTok{,} \VariableTok{normpdf}\NormalTok{(}\VariableTok{x}\OperatorTok{,} \FloatTok{1}\OperatorTok{,} \VariableTok{sqrt}\NormalTok{(}\FloatTok{2}\OperatorTok{/}\VariableTok{ns}\NormalTok{(}\VariableTok{i}\NormalTok{)))}\OperatorTok{,} \SpecialStringTok{\textquotesingle{}r\textquotesingle{}}\OperatorTok{,} \SpecialStringTok{\textquotesingle{}LineWidth\textquotesingle{}}\OperatorTok{,} \FloatTok{2}\NormalTok{)}\OperatorTok{;}
    \VariableTok{title}\NormalTok{(}\VariableTok{sprintf}\NormalTok{(}\SpecialStringTok{\textquotesingle{}n=\%d\textquotesingle{}}\OperatorTok{,} \VariableTok{ns}\NormalTok{(}\VariableTok{i}\NormalTok{)))}\OperatorTok{;}
\KeywordTok{end}
\VariableTok{sgtitle}\NormalTok{(}\SpecialStringTok{\textquotesingle{}CLT for χ²(1): Mean→Gaussian\textquotesingle{}}\NormalTok{)}\OperatorTok{;}
\end{Highlighting}
\end{Shaded}

\emph{Next Module: {Module 12 — Introduction to Statistics}}

\end{PSbox}
